  \recommended \item \label{exer:PowersOfDiags}
    对角矩阵的幂具有什么形式？
    \begin{answer}
      对于任意整数 \( p \)，我们有下式。
      \begin{equation*}
        \begin{mat}
          d_1  &0      &   \\
          0    &\ddots &   \\
               &       &d_n
        \end{mat}^p=
        \begin{mat}
          d_1^p  &0      &   \\
          0      &\ddots &   \\
                 &       &d_n^p
        \end{mat}
     \end{equation*} 
    \end{answer}
  \item 
     给出两个同阶但不相似的对角矩阵。
     任意两个不同的对角矩阵是否必定来自不同的相似类？
     \begin{answer}
       这两个矩阵不相似 
       \begin{equation*}
          \begin{mat}[r]
             0  &0  \\
             0  &0
          \end{mat}
          \qquad
          \begin{mat}[r]
             1  &0  \\
             0  &1
          \end{mat}
       \end{equation*}
       因为每一个在各自的相似类中都是单独的。

       对于后半问，这两个
       \begin{equation*}
          \begin{mat}[r]
             2  &0  \\
             0  &3
          \end{mat}
          \qquad
          \begin{mat}[r]
             3  &0  \\
             0  &2
          \end{mat}
       \end{equation*}
       通过把基从
       \( \sequence{\vec{\beta}_1,\vec{\beta}_2} \) 变为
       \( \sequence{\vec{\beta}_2,\vec{\beta}_1} \) 的矩阵而相似。
       （\textit{问题。}
        两个对角矩阵相似当且仅当它们的对角元互为排列？）  
    \end{answer}
  \item 
    给出一个非奇异的对角矩阵。
    对角矩阵有可能是奇异的吗？
    \begin{answer}
      对比这两个矩阵。
      \begin{equation*}
         \begin{mat}[r]
           2  &0  \\
           0  &1
         \end{mat}
         \qquad
         \begin{mat}[r]
           2  &0  \\
           0  &0
         \end{mat}
      \end{equation*}
      第一个是非奇异的，第二个是奇异的。  
     \end{answer}
  \recommended \item
    证明：若对角线上没有元素为零，则对角矩阵的逆就是由各元素的
    逆构成的对角矩阵。
    当对角元为零时会发生什么？
    \begin{answer}  
       要验证对角矩阵的逆是由各元素的逆构成的对角矩阵，只需直接相乘。
       \begin{equation*}
          \begin{mat}
             a_{1,1}  &0                \\
             0        &a_{2,2}          \\
                      &       &\ddots    \\
                      &       &      &a_{n,n}
          \end{mat}
          \begin{mat}
            1/a_{1,1}  &0                \\
             0        &1/a_{2,2}          \\
                      &       &\ddots    \\
                      &       &      &1/a_{n,n}
          \end{mat}
      \end{equation*}
      （证明它是左逆同样容易。）

      若某个对角元为零，则该对角矩阵是奇异的；它的行列式为零。  
    \end{answer}
  \item 
    以 \nearbyexample{ex:DiagUpperTrian} 结尾的方程
    \begin{equation*}
       \begin{mat}[r]
         1  &1  \\
         0  &-1
       \end{mat}^{-1}
       \begin{mat}[r]
         3  &2  \\
         0  &1
       \end{mat}
       \begin{mat}[r]
         1  &1  \\
         0  &-1
       \end{mat}
       =
       \begin{mat}[r]
         3  &0  \\
         0  &1
       \end{mat}
    \end{equation*}
    看起来有点别扭，因为 $P$ 必须取第一个矩阵，
    而它是以逆的形式出现的；$P^{-1}$ 则取第一个矩阵的逆，
    于是两个 $-1$ 次幂相消，这个矩阵出现时
    就不带上标 $-1$ 了。
    \begin{exparts}
      \partsitem 验证下面这个形式更好看的方程成立。
        \begin{equation*}
           \begin{mat}[r]
             3  &0  \\
             0  &1
           \end{mat}
           =
           \begin{mat}[r]
             1  &1  \\
             0  &-1
           \end{mat}
           \begin{mat}[r]
             3  &2  \\
             0  &1
           \end{mat}
           \begin{mat}[r]
             1  &1  \\
             0  &-1
           \end{mat}^{-1}
        \end{equation*}
      \partsitem 上一小问是巧合吗？
        还是我们总能交换 $P$ 与 $P^{-1}$？
   \end{exparts}
   \begin{answer}
     \begin{exparts}
       \partsitem 验证很容易。
         \begin{equation*}
           \begin{mat}[r]
             1  &1  \\
             0  &-1
           \end{mat}
           \begin{mat}[r]
             3  &2  \\
             0  &1
           \end{mat}
           =
           \begin{mat}[r]
             3  &3  \\
             0  &-1
           \end{mat}
           \qquad
           \begin{mat}[r]
             3  &3  \\
             0  &-1
           \end{mat}
           \begin{mat}[r]
             1  &1  \\
             0  &-1
           \end{mat}^{-1}
           =
           \begin{mat}[r]
             3  &0  \\
             0  &1
           \end{mat}
         \end{equation*}
        \partsitem 这是一个巧合，也就是说，若 $T=PSP^{-1}$，
          则 $T$ 不必等于 $P^{-1}SP$。
          甚至对于对角矩阵~$D$，条件 $D=PTP^{-1}$
          也不能推出 $D$ 等于 $P^{-1}TP$。
          来自 \nearbyexample{ex:DiagTwoByTwo} 的矩阵说明了这一点。
          \begin{equation*}
            \begin{mat}[r]
              1  &2  \\
              1  &1
            \end{mat}
            \begin{mat}[r]
              4  &-2  \\
              1  &1
            \end{mat}
            =
            \begin{mat}[r]
              6  &0  \\
              5  &-1
            \end{mat}
            \qquad
            \begin{mat}[r]
              6  &0  \\
              5  &-1
            \end{mat}
            \begin{mat}[r]
              1  &2  \\
              1  &1
            \end{mat}^{-1}
            =
            \begin{mat}[r]
              -6  &12  \\
              -6  &11
            \end{mat}
          \end{equation*}
     \end{exparts}
   \end{answer}
  \item 
    证明：在 \nearbyexample{ex:DiagUpperTrian}
    中用于对角化的 $P$ 不是唯一的。
    \begin{answer}
      该矩阵的列就是与各个 $x$ 相联系的向量。
      具体的选取，以及选取的顺序，都是任意的。
      例如，交换这两列就能得到另一个矩阵。
    \end{answer}
  \item 
    求该矩阵的幂的一个公式。
    \textit{提示}：~参见 \nearbyexercise{exer:PowersOfDiags}。
    \begin{equation*}
      \begin{mat}[r]
        -3  &1  \\
        -4  &2
      \end{mat}
    \end{equation*}
    \begin{answer}
      先对角化，再取对角矩阵的幂，可知
      \begin{equation*}
        \begin{mat}[r]
          -3  &1  \\
          -4  &2
        \end{mat}^k
        =
        \frac{1}{3}
        \begin{mat}[r]
          -1  &1  \\
          -4  &4
        \end{mat}
        +(\frac{-2}{3})^k
        \begin{mat}[r]
           4  &-1 \\
           4  &-1
        \end{mat}.
      \end{equation*}  
     \end{answer}
  \item 
    我们可以问一问对角化与矩阵运算如何相互作用。
    设 \( \map{t,s}{V}{V} \) 都是可对角化的。
    对所有标量 \( c \)，\( ct \) 都可对角化吗？
    那 \( t+s \) 呢？
    \( \composed{t}{s} \) 呢？
    \begin{answer}
      是的，由本小节最后一个定理可知 \( ct \) 可对角化。

      不是，\( t+s \) 不一定可对角化。
      直观地说，当这两个映射关于不同的基对角化时
      （即当它们不是
      \definend{可同时对角化的}）问题就出现了。
      具体地，这两个矩阵都可对角化，但它们的和不可对角化：
      \begin{equation*}
         \begin{mat}[r]
            1  &1  \\
            0  &0
         \end{mat}
         \qquad
         \begin{mat}[r]
           -1  &0  \\
            0  &0
         \end{mat}
      \end{equation*}
      （第二个已经是对角矩阵；第一个
      参见 \nearbyexercise{exer:DiagTheseA}）。
      其和不可对角化，因为它的平方是零矩阵。 

     同样的直觉提示 \( \composed{t}{s} \) 也不
     可对角化。
     这两个矩阵都可对角化，但它们的乘积不可对角化：
     \begin{equation*}
        \begin{mat}[r]
           1  &0  \\
           0  &0
        \end{mat}
        \qquad
        \begin{mat}[r]
           0  &1  \\
           1  &0
        \end{mat}
     \end{equation*}
     （第二个参见 \nearbyexercise{exer:DiagTheseB}）。   
    \end{answer}
  \item 
    证明：这种形式的矩阵不可对角化。
    \begin{equation*}
       \begin{mat}[r]
          1  &c  \\
          0  &1
       \end{mat}
       \qquad c\neq 0
    \end{equation*}
    \begin{answer}
      若
      \begin{equation*}
         P
         \begin{mat}
            1  &c  \\
            0  &1
         \end{mat}
         P^{-1}
         =
         \begin{mat}
            a  &0  \\
            0  &b
         \end{mat}
      \end{equation*}
      则
      \begin{equation*}
         P
         \begin{mat}
            1  &c  \\
            0  &1
         \end{mat}
         =
         \begin{mat}
            a  &0  \\
            0  &b
         \end{mat}
         P
      \end{equation*}
      于是
      \begin{align*}
         \begin{mat}
            p  &q  \\
            r  &s
         \end{mat}
         \begin{mat}
            1  &c  \\
            0  &1
         \end{mat}
         &=
         \begin{mat}
            a  &0  \\
            0  &b
         \end{mat}
         \begin{mat}
            p  &q  \\
            r  &s
         \end{mat}        \\
         \begin{mat}
            p  &cp+q  \\
            r  &cr+s
         \end{mat}
         &=
         \begin{mat}
            ap  &aq  \\
            br  &bs
         \end{mat}
      \end{align*}
      \( 1,1 \) 元表明 \( a=1 \)，于是 \( 1,2 \) 元
      表明 \( pc=0 \)。
      由于 \( c\neq 0 \)，这意味着 \( p=0 \)。
      \( 2,1 \) 元表明 
      \( b=1 \)，于是 \( 2,2 \) 元表明
      \( rc=0 \)。
      由于 \( c\neq 0 \)，这意味着 \( r=0 \)。
      但若 \( p \) 与 \( r \) 都是 \( 0 \)，则 \( P \) 不
      可逆。  
     \end{answer}
  \item 
    证明：下面每一个矩阵都可对角化。
    \begin{exparts*}
      \partsitem
       \( \begin{mat}[r]
             1  &2  \\
             2  &1
          \end{mat}  \)
      \partsitem
       \( \begin{mat}
             x  &y  \\
             y  &z
          \end{mat}
          \qquad \text{$x,y,z$ 为标量}  \)
    \end{exparts*}
    \begin{answer}
      \begin{exparts}
      \partsitem 使用 $\nbyn{2}$
        矩阵的逆的公式可得下式。
        \begin{multline*}
           \begin{mat}
              a  &b  \\
              c  &d
           \end{mat}
           \begin{mat}[r]
              1  &2  \\
              2  &1
           \end{mat}
           \cdot\frac{1}{ad-bc}\cdot
           \begin{mat}
              d  &-b \\
             -c  &a
           \end{mat}                                  \\
           =\frac{1}{ad-bc}
           \begin{mat}
              ad+2bd-2ac-bc    &-ab-2b^2+2a^2+ab \\
              cd+2d^2-2c^2-cd  &-bc-2bd+2ac+ad
           \end{mat}
        \end{multline*}
        现在选取标量 \( a,\ldots,d \)，使得
        \( ad-bc\neq 0 \) 且 \( 2d^2-2c^2=0 \) 且 \( 2a^2-2b^2=0 \)。
        例如，下面这组就行。
        \begin{equation*}
          \begin{mat}[r]
            1  &1  \\
            1  &-1
          \end{mat}
          \begin{mat}[r]
            1  &2  \\
            2  &1
          \end{mat}
          \cdot\frac{1}{-2}\cdot
          \begin{mat}[r]
            -1  &-1  \\
            -1  &1
          \end{mat}
          =
          \frac{1}{-2}
          \begin{mat}[r]
            -6  &0   \\
             0  &2
          \end{mat}
        \end{equation*}
      \partsitem 同上，
        \begin{multline*}
           \begin{mat}
              a  &b  \\
              c  &d
           \end{mat}
           \begin{mat}
              x  &y  \\
              y  &z
           \end{mat}
           \cdot\frac{1}{ad-bc}\cdot
           \begin{mat}
              d  &-b \\
             -c  &a
           \end{mat}                               \\
           =\frac{1}{ad-bc}
           \begin{mat}
              adx+bdy-acy-bcz    &-abx-b^2y+a^2y+abz \\
              cdx+d^2y-c^2y-cdz  &-bcx-bdy+acy+adz
           \end{mat}
        \end{multline*}
        我们要找标量 \( a,\ldots,d \)，使得
        \( ad-bc\neq 0 \) 且
        \( -abx-b^2y+a^2y+abz=0 \)
        且 \( cdx+d^2y-c^2y-cdz=0 \)，无论
        \( x \)、\( y \)、\( z \)
        取什么值。

        首先，我们假定 \( y\neq 0 \)，否则所给矩阵已经
        是对角矩阵。
        我们将使用该假定，因为若我们（随意地）令
        \( a=1 \) 则得到
        \begin{align*}
           -bx-b^2y+y+bz
           &=0              \\
           (-y)b^2+(z-x)b+y
           &=0
        \end{align*}
        而求根公式给出
        \begin{equation*}
           b=\frac{-(z-x)\pm\sqrt{(z-x)^2-4(-y)(y)} }{-2y}
           \qquad
           y\neq 0
        \end{equation*}
        （注意若 \( x \)、\( y \)、\( z \) 为实数，则这两个
         \( b \) 都是实数，因为判别式为正）。
        同理，若我们（随意地）令 \( c=1 \) 则
        \begin{align*}
           dx+d^2y-y-dz
           &=0              \\
           (y)d^2+(x-z)d-y
           &=0
        \end{align*}
        在这里我们得到
        \begin{equation*}
           d=\frac{-(x-z)\pm\sqrt{(x-z)^2-4(y)(-y)} }{2y}
           \qquad
           y\neq 0
        \end{equation*}
        （同上，若 \( x,y,z\in\Re \)，则此判别式为正，
        故实对称的 \( \nbyn{2} \) 矩阵相似于一个实的
        对角矩阵）。

        作为检验，我们取 \( x=1 \)、\( y=2 \)、\( z=1 \)。
        \begin{equation*}
           b=\frac{0\pm\sqrt{0+16} }{-4}=\mp 1
           \qquad
           d=\frac{0\pm\sqrt{0+16} }{4}=\pm 1
        \end{equation*}
        注意并非所有四个选择 \( (b,d)=(+1,+1),\dots,(-1,-1) \)
        都满足 \( ad-bc\neq 0 \)。
      \end{exparts} 
    \end{answer}
\index{diagonalizable|)}
\end{exercises}
